\section{Lean 4: equivalence classes form a partition}
\label{sec:lean-bridge-05}
Let $r$ be a reflexive, symmetric, and transitive relation on a
type $\alpha$. We define the class of $a$ by
$[a]=\{b:r(b,a)\}$. Reflexivity proves that every $a$ lies
in its own class, so the classes cover the carrier.

If $x$ lies in both $[a]$ and $[b]$, symmetry followed by
transitivity gives $a\,r\,b$. Consequently every element of
$[a]$ lies in $[b]$, and conversely. Thus intersecting
classes are equal. This establishes the two characteristic
partition properties without assuming that the relation is
equality.

\lstinputlisting[language=LeanFour,basicstyle=\ttfamily\footnotesize,
breaklines=true,columns=fullflexible,keepspaces=true,frame=single,
firstline=73,lastline=97]
{../../formal/lean/Con2physFormal/Bridges.lean}
